\bookchapter{Kapitel 2: Räum dein digitales Zimmer auf}

\chapterintro{Mit der Zeit wird ein Smartphone unordentlich — gute Gewohnheiten beginnen mit der Dateiverwaltung. Wie dein Zimmer regelmäßig aufgeräumt werden muss, braucht auch dein digitaler Raum Pflege. Ein aufgeräumtes Smartphone findet Dinge nicht nur schneller, sondern läuft auch flüssiger.}

\booksection{Eins. Das dreistufige Speichersystem: Dein digitales Lager}

\begin{enumerate}
  \item \textbf{Interner Speicher}: Der Kernbereich – schnell, hier wohnen System und Apps. Er ist der kostbarste Platz deines Handys: klug einteilen.
  \item \textbf{SD-Karte (falls unterstützt)}: Die erweiterte Garderobe, etwas langsamer, geeignet für Filme, Fotos und andere große Dateien. Achtung: Viele Flagship-Smartphones unterstützen keine SD-Karten mehr.
  \item \textbf{Cloud-Speicher}: Tresor plus Sync-Cloud – schützt vor Verlust und synchronisiert über mehrere Geräte. iCloud, Baidu Netdisk und OneDrive sind gängige Optionen.
\end{enumerate}

\booksection{Zwei. Lerne den Dateimanager kennen}

Zwei. Lerne den Dateimanager kennen

\booksubsection{Android hat meist eine [Dateien]- oder [Dateimanager]-App, iOS die integrierte [Dateien]-App. Sie sind dein „Explorer“ für alle Dateien auf dem Smartphone.}

\begin{itemize}
  \item \textbf{Bedeutung gängiger Ordner:}DCIM
  \item \textbf{: Fotos und Videos der Kamera.}: Andere Bilder, wie Screenshots und gespeicherte Bilder.
  \item \textbf{Download}: Vom Browser heruntergeladene Dateien.
  \item \textbf{Music}: Musikdateien.
  \item \textbf{Movies}: Videodateien.
  \item \textbf{Documents}: Dokumente.
  \item \textbf{Android}: Der private Bereich einer App, für andere Apps nicht zugänglich. Wird beim Deinstallieren der App gelöscht.
\end{itemize}

\booksection{Der Aufbau des Dateisystems: ein tieferer Blick}

Android basiert auf Linux und hat eine gewisse Hierarchie im Dateisystem. Auch wenn normale Nutzer das nicht im Detail verstehen müssen, helfen einige Grundbegriffe.

\booksubsection{Speicherpartitionen:}

\begin{itemize}
  \item \textbf{/system}: Die Systempartition mit den Dateien des Betriebssystems. Für normale Nutzer gesperrt, Zugriff nur mit Root-Rechten.
  \item \textbf{/data}: Die Datenpartition für Nutzer- und App-Daten. Das ist dein alltäglicher Speicherplatz.
  \item \textbf{/sdcard oder /storage/emulated/0}: Der nutzerzugängliche Speicher – hier liegen deine Fotos und Downloads.
\end{itemize}

\booksubsection{: Dateien, die eine App im nutzerzugänglichen Bereich ablegt, wie Fotos oder Downloads. Werden beim Deinstallieren nicht zwingend gelöscht.}

\begin{itemize}
  \item \textbf{Interner Speicher}: Der private Bereich einer App, für andere Apps unzugänglich. Wird beim Deinstallieren gelöscht.
  \item \textbf{Externer Speicher}: Dateien, die eine App im nutzerzugänglichen Bereich ablegt, etwa Fotos und Downloads. Beim Deinstallieren nicht zwangsläufig gelöscht.
\end{itemize}

\booksection{Backup und Wiederherstellung: eine Versicherung für deine Daten}

Daten sind unbezahlbar – Backups sind eine Gewohnheit, die du dir angewöhnen musst.

\booksubsection{Backup-Strategien:}

\begin{enumerate}
  \item \textbf{Cloud-Backup}: Automatisch und bequem, braucht aber Netz und Cloud-Speicher. iCloud, Huawei Cloud, Xiaomi Cloud und andere bieten das.
  \item \textbf{Lokales Backup}: Per Kabel mit dem Computer verbinden und mit dem offiziellen Tool sichern. Schnell und raumsparend, aber manueller Aufwand.
  \item \textbf{Doppeltes Backup}: Cloud- plus lokales Backup – am sichersten.
\end{enumerate}

\booksubsection{Was sichert man?}

\begin{itemize}
  \item \textbf{Fotos und Videos}: Platzrauber und zugleich das Kostbarste.
  \item \textbf{Kontakte}: Am einfachsten im Cloud-Konto aufgehoben.
  \item \textbf{WeChat-Chatverläufe}: Über die integrierte Backup-Funktion von WeChat sichern.
  \item \textbf{Wichtige Dokumente}: In der Cloud oder auf dem Computer ablegen.
\end{itemize}

\booksection{Dateisicherheit: schütze deine Privatsphäre}

\begin{itemize}
  \item \textbf{Verschlüsselung sensibler Dateien}: Viele Handys haben einen „sicheren Ordner“ oder „privaten Raum“: sensible Dateien hinein, Zugriff nur mit Passwort oder Fingerabdruck.
  \item \textbf{Löschen ist nicht Vernichten}: Normales Löschen markiert den Speicher nur als überschreibbar – die Daten bleiben. Zum echten Entfernen brauchst du ein sicheres Löschtool oder mehrfaches Überschreiben.
  \item \textbf{Cloud-Sicherheit}: Seriösen Cloud-Anbieter wählen, Zwei-Faktor-Authentifizierung aktivieren, angemeldete Geräte regelmäßig prüfen.
\end{itemize}

\booksection{Speicher knapp? Zeit für eine Großreinigung}

\begin{itemize}
  \item \textbf{Cache leeren}: Das Kurzzeitgedächtnis der Apps – Leeren schafft Platz, unter Einstellungen > Speicher.
  \item \textbf{Große Dateien finden}: Im Dateimanager nach Größe sortieren und die „Sündenböcke“ entlarven.
  \item \textbf{Rückstände entfernen}: Nach der Deinstallation bleiben manchmal Mülldateien zurück – von Hand löschen.
  \item \textbf{Galerie aufräumen}: Doppelt fotografierte, unscharfe und unnötige Screenshots löschen.
  \item \textbf{WeChat aufräumen}: WeChat ist ein Speicherriese – Chatverläufe und Cache regelmäßig zu leeren lohnt sich.
\end{itemize}

\booksection{Alle gängigen Dateiformate erklärt}

Die Dateiendungen auf dem Handy können verwirren – wer sie versteht, verwaltet seine Dateien besser:

\begin{center}
\begin{tabularx}{\linewidth}{|X|X|X|X|}
\hline
\textbf{Format} & \textbf{Typ} & \textbf{Eigenschaften} & \textbf{Empfohlene Öffnungsweise} \tabularnewline
\hline
JPG/JPEG & Bild & Verlustbehaftete Kompression, geringe Größe, beste Kompatibilität. Standardformat für Smartphone-Fotos. & Album / beliebiger Bildbetrachter \tabularnewline
PNG & Bild & Verlustfreie Kompression, unterstützt transparente Hintergründe. Häufig für Screenshots und Icons. Größer als JPG. & Album / beliebiger Bildbetrachter \tabularnewline
MP3 & Audio & Verlustbehaftete Kompression, geringe Größe, hervorragende Kompatibilität. Das universellste Audioformat. & Beliebiger Musikplayer \tabularnewline
FLAC & Audio & Verlustfreie Kompression, Klangqualität wie CD, etwa 3-5 mal so groß wie MP3. Erste Wahl für Audiophile. & Musik-Player / Profi-Player-App \tabularnewline
WAV & Audio & : Unkomprimiertes Original-Audio, beste Qualität, riesige Dateien. Üblich in der Profi-Aufnahme. & Musik-Player \tabularnewline
M4A & Audio & AAC-Codierung, bei gleicher Bitrate bessere Qualität als MP3. Im Apple-Ökosystem verbreitet. & Apple Music / Musikplayer \tabularnewline
MP4 & Video & Das universellste Video-Containerformat, hervorragende Kompatibilität. Standardformat für Smartphone-Videos. & Album / beliebiger Videoplayer \tabularnewline
MKV & Video & Offenes Containerformat, unterstützt mehrere Tonspuren und Untertitel. Häufig bei Filmressourcen. & VLC / nPlayer \tabularnewline
Word-Dokument. DOCX ist ein komprimiertes Format mit geringerer Größe. DOC ist das alte Format. & Dokumente & : Word-Dokument. DOCX ist das komprimierte Format und kleiner; DOC das alte Format. & WPS / Office \tabularnewline
PPT/PPTX & Präsentationen & PPT/PPTX & WPS / Office \tabularnewline
Excel-Tabelle. XLSX ist ein komprimiertes Format. & Tabellen & : Excel-Tabelle; XLSX ist das komprimierte Format. & WPS / Office \tabularnewline
PDF & Dokumente & : Plattformübergreifendes Format mit festem, nicht editierbarem Layout. Üblich für Verträge, Arbeiten und E-Books. & Adobe Reader / WPS \tabularnewline
ZIP & Komprimiert & Optisches Disc-Image, komplette Kopie des Disc-Inhalts. Häufig bei System-Installationspaketen. & Systemeigene Funktion / Entpack-App \tabularnewline
RAR & Komprimiert & Optisches Disc-Image, komplette Kopie des Disc-Inhalts. Häufig bei System-Installationspaketen. & Entpack-App \tabularnewline
ISO & Images & : Disc-Image, eine komplette Kopie einer CD/DVD. Üblich für System-Installationspakete. & Entpacker-App / virtuelles Laufwerk \tabularnewline
HTML & Webseiten & : Webseitendatei, im Browser öffnen. Üblich für gespeicherte Webseiten. & Browser \tabularnewline
\hline
\end{tabularx}
\end{center}

\begin{tcbpractice}{Übung: räum den Ordner „Downloads“ auf}
Dateimanager öffnen > in den Ordner „Downloads“ gehen. Was liegt da drin? Unnötiges löschen, Nützliches sortieren (z. B. neue Ordner „Arbeitsunterlagen“ und „Sticker“ anlegen). Gleichzeitig einen Blick auf den Speicher werfen: Welche Dateitypen belegen am meisten?

\end{tcbpractice}

\booksection{Foto- und Dateiverwaltung: ordentliche Vermögensverwaltung}

Im letzten Abschnitt hast du „Ordner verstehen“ gelernt – aber verstehen allein reicht nicht. In diesem Abschnitt geht es darum,\textbf{wie du eine dauerhaft tragfähige Verwaltungsregel aufbaust}, damit deine Dateien jederzeit auffindbar, nicht verloren und nicht ausspioniert sind.

\booksubsection{Dateiklassifizierung und Benennungsregeln}

Das Chaos entsteht aus Regellosigkeit. Eine Regel aufzustellen ist einfach – entscheidend ist\textbf{dranzubleiben}.

\begin{itemize}
  \item \textbf{Benennung nach „Zeit + Ereignis“}: Zeitstempel und Ereignisbeschreibung im Ordnernamen helfen später bei Erinnerung und Suche. Zum Beispiel\texttt{202606\_Reisebelege},\texttt{2025\_Jahresrückblick},\texttt{2024\_Kindgeburtstag}. So siehst du einerseits die Zeitachse und weißt andererseits auf einen Blick, worum es geht.
  \item \textbf{Dreistufiges Klassifikationssystem}:
  \begin{itemize}
    \item \textbf{Stufe 1}: Nach Dateityp (Dokument / Bild / Video / Audio)
    \item \textbf{Stufe 2}: Nach Nutzungsszenario (Arbeit / Leben / Lernen)
    \item \textbf{.}: Nach konkrettem Ereignis oder Projekt (z. B. „2026Q2\_ProjektA")
  \end{itemize}
  \item \textbf{Monatlicher Aufräumplan}: Lege jeden Monat einen festen Tag fest (am besten den Gehaltstag oder den 1. des Monats), um das „Download"-Verzeichnis konzentriert aufzuräumen:\textbf{,}. Zwölf kleine Aufräumaktionen im Jahr sind viel entspannter als eine Jahresend-Generalreinigung.
\end{itemize}

\begin{tcbpractice}{,}
Such dir ein freies Wochenende aus und gehe in dieser Reihenfolge vor:

\begin{enumerate}
  \item Lege im Dateimanager einen Arbeitsordner für diesen Monat im Format\texttt{2026Q2\_ProjektA}:
  \item Gehe alle Dateien dieser Woche im „Download"-Verzeichnis durch: Nützliches behalten und archivieren, Nutzloses löschen.
  \item Erstelle für die Familienfotos dieses Jahr Ordner wie\texttt{2025\_Familie}oder\texttt{2026\_Familie}.
\end{enumerate}

\end{tcbpractice}

\booksubsection{Strategien für redundante Dateien}

Mit der Zeit wird ein Smartphone „dicker“ — meist sind redundante Dateien schuld.

\begin{itemize}
  \item \textbf{Automatische 72-Stunden-Aufräumung}: Temporäre Dateien, die über soziale Apps wie WeChat oder QQ empfangen werden (Sticker, Chat-Bilder, Kurzvideos), sind nach einmal Ansehen meist nutzlos. Unter [WeChat] [Einstellungen] [Allgemein] [Speicherplatz] die „automatische Bereinigung" aktivieren und die Aufbewahrung auf 72 Stunden setzen. Das Smartphone ist sofort aufgeräumter, ohne wichtige Aufzeichnungen zu verlieren.
  \item \textbf{Gezielte Speicherbereinigung}: Android [Einstellungen] > [Speicher] > [Reinigen \& beschleunigen]; iOS [Einstellungen] > [Allgemein] > [iPhone-Speicher] – das System listet nach „Cache / temporäre Dateien / entfernbare Apps / große Dateien“ auf. Prioritäten: (1) Cache-Dateien; (2) abgelaufene Installationspakete (alte APK/IPA nach App-Updates); (3) doppelte Dateien (Fotos mit gleichem Namen und Größe).
  \item \textbf{Regeln für die automatische Sortierung}: Die Galerie-Apps der meisten Marken können intelligent sortieren: Personen, Landschaft, Dokumente, Screenshots. Einmal aktiviert, landen neue Gruppenfotos automatisch unter „Personen > Name“ und gescannte Verträge unter „Dokumente“.\textbf{Mit der monatlichen Aufräumroutine kombiniert, sinkt die Chance auf „Ich find’s nicht“ enorm}.
\end{itemize}

\booksubsection{Regeln für den sicheren Umgang mit Bildern}

Bilder sind der „sensibelste“ Datentyp im Handy – ein einziger Schnappschuss kann Ort, Gerät und Zeit verraten. Wir empfehlen einen\textbf{Workflow mit „doppeltem Schutz“}:

\begin{tcbwarn}{Standardablauf für Ausweis-/Sensible-Screenshots}
\begin{enumerate}
  \item \textbf{Sichtbare Informationen verpixeln}: Mit systemeigenen oder Drittanbieter-Tools (wie Meitu XiuXiu, Snapseed)\textbf{die mittleren Ziffern der Personalausweisnummer, das Kennzeichen, die Hausnummer, den Namen und Beträge}schwärzen.
  \item \textbf{Metadaten entfernen}: Beim Teilen im Album „Teilen" oder „Exportieren" wählen — das System fragt, ob Standortinformationen erhalten bleiben sollen,\textbf{„Nicht erhalten"}wählen. So werden Aufnahmegerät, GPS-Koordinaten und Aufnahmezeit nicht mit dem Bild verschickt.
  \item \textbf{Temporäres Wasserzeichen}(optional): Auf sensiblen Screenshots ein Textwasserzeichen wie „Nur zur Verwendung für XX" anbringen, um Missbrauch zu vermeiden.
\end{enumerate}

\end{tcbwarn}

.

\begin{itemize}
  \item \textbf{Systemeigen}: Unter iOS [Einstellungen] [Datenschutz \& Sicherheit] [Kamera] lässt sich „Genaue Standortbestimmung" deaktivieren; beim Teilen mit „Nur aktuelles Foto" werden\textbf{die meisten EXIF-Daten automatisch entfernt}:
  \item \textbf{Android}: Die Album-Apps der meisten Marken fragen beim „Teilen", ob die Standortinformationen entfernt werden sollen — einfach anhaken.
  \item \textbf{Professionelle Tools}: ExifTool (Desktop), Metapho (iOS), Photo Metadata Remover (Android) können alle Metadaten restlos löschen.
\end{itemize}

Für Bilder, die\textbf{langfristig aufbewahrt}werden müssen (Ausweis, Bankkarte, Einwohnermeldeamt-Scan, Reisepass), empfehlen wir die Ablage in einem\textbf{verschlüsselten Album}.

\begin{itemize}
  \item iPhone: [Fotos] Foto auswählen Ausblenden (nur nach Face-ID-Verifizierung sichtbar).
  \item Huawei / Xiaomi: [Album] Foto lange drücken „In Tresor verschieben".
  \item Haeufige Methoden zum Entfernen von EXIF-Metadaten (Grundoperationen, tiefergehender Schutz in Kapitel 6 „Fotos anonymisieren"):\textbf{360 Safe Album},\textbf{Tencent Album Manager (Privatalbum)}.\textbf{Android}u. a. mit separatem Passwortschutz und besserer Tarnung.
\end{itemize}
